\documentclass[10pt,a4paper]{article}
\usepackage[margin=18mm,headheight=14pt]{geometry}
\usepackage{fontspec}
\setmainfont{lmroman10-regular.otf}[BoldFont=lmroman10-bold.otf,ItalicFont=lmroman10-italic.otf,BoldItalicFont=lmroman10-bolditalic.otf]
\setsansfont{lmsans10-regular.otf}[BoldFont=lmsans10-bold.otf]
\setmonofont{lmmonolt10-regular.otf}[BoldFont=lmmonolt10-bold.otf,ItalicFont=lmmonolt10-oblique.otf]
\usepackage[ngerman]{babel}
\usepackage{amsmath,amssymb,booktabs,tabularx,array,fancyhdr,listings,xcolor}
\usepackage[hidelinks]{hyperref}
\pagestyle{fancy}\fancyhf{}
\fancyhead[L]{\sffamily AI-Hardware · Tag 5}
\fancyhead[R]{\sffamily Logikgatter und Boolesche Algebra}
\fancyfoot[L]{\small Nachschlagen · Verstehen · Schaltungen prüfen}
\fancyfoot[R]{\thepage\ / 8}
\setlength{\parindent}{0pt}\setlength{\parskip}{5pt}
\setlength{\tabcolsep}{7pt}\renewcommand{\arraystretch}{1.18}
\lstset{language=Python,basicstyle=\ttfamily\small,
 keywordstyle=\bfseries,commentstyle=\itshape\color{black},
 stringstyle=\ttfamily\color{black},showstringspaces=false,
 frame=single,rulecolor=\color{black},framesep=5pt,
 columns=fullflexible,keepspaces=true,breaklines=true,tabsize=4,
 aboveskip=8pt,belowskip=8pt}
\newcommand{\titlepagepart}[2]{\vspace*{-4mm}{\sffamily\small #1}\par{\sffamily\LARGE\bfseries #2}\par\vspace{2mm}}
\newcommand{\block}[1]{\par\vspace{2mm}{\sffamily\large\bfseries #1}\par}
\newcommand{\code}[1]{\texttt{#1}}

\usepackage{tikz}
\usetikzlibrary{shapes.gates.logic.US,positioning,calc,patterns}
\tikzset{gate/.style={draw,minimum width=15mm,minimum height=10mm,font=\sffamily\small},wire/.style={thick}}
\newcommand{\source}[1]{\par{\footnotesize\itshape Buchbezug: #1. Abbildungen in diesem Skript sind eigene Zeichnungen.}}
\begin{document}
\titlepagepart{01 / DER EINSTIEG}{Was macht ein Logikgatter?}
\textbf{Dein Tagesziel:} Logische Bedingungen als Gleichung, Wahrheitstabelle
und Schaltung darstellen. NOT, AND, OR und XOR aus NAND bauen; De Morgan
nachweisen. Außerdem verstehen, warum reale Spannungen zuverlässig als 0 oder 1 gelten.

Ein \textbf{Logikgatter} verarbeitet Eingangssignale und erzeugt einen Ausgang.
Heute bedeutet 1 „wahr“ und 0 „falsch“. Anders als an Tag 4 addieren wir hier
nicht automatisch Stellenwerte: Wir untersuchen zunächst einzelne logische Signale.

\block{Ein Beispiel: Zwei Bedingungen müssen erfüllt sein}
Eine vereinfachte Maschinenfreigabe soll nur aktiv sein, wenn der Startknopf
gedrückt \emph{und} die Klappe geschlossen ist. $A$ steht für Start,
$B$ für Klappe geschlossen, $Y$ für Freigabe. Das ist AND.
\begin{center}
\begin{tikzpicture}
\node[and gate US,draw,logic gate inputs=nn,scale=1.6] (g) at (0,0) {};
\draw[wire] (g.input 1)--++(-1.2,0) node[left] {$A$};
\draw[wire] (g.input 2)--++(-1.2,0) node[left] {$B$};
\draw[wire] (g.output)--++(1.5,0) node[right] {$Y=A\cdot B$};
\end{tikzpicture}
\end{center}
\begin{center}
\begin{tabular}{ccl}\toprule
$A$&$B$&$Y$: Freigabe\\\midrule
0&0&0: keine Bedingung erfüllt\\
0&1&0: Start fehlt\\
1&0&0: Klappe ist offen\\
1&1&1: beide Bedingungen erfüllt\\\bottomrule
\end{tabular}
\end{center}
Dies ist eine \textbf{Wahrheitstabelle}. Sie enthält für zwei Eingänge alle
$2^2=4$ Kombinationen. Für drei Eingänge braucht sie $2^3=8$ Zeilen.
Die Eingänge werden wie Binärzahlen aufgezählt: 00, 01, 10, 11.
Das Beispiel erklärt eine logische Funktion, keine vollständige Sicherheitssteuerung.

\block{Vier Darstellungen derselben Funktion}
\begin{center}
\begin{tikzpicture}[node distance=12mm]
\node[gate] (a) {Worte};
\node[gate,right=of a] (b) {Gleichung};
\node[gate,right=of b] (c) {Gatter};
\node[gate,right=of c] (d) {Tabelle};
\draw[<->] (a)--(b);\draw[<->] (b)--(c);\draw[<->] (c)--(d);
\end{tikzpicture}
\end{center}
Worte sagen, \emph{was} passieren soll. Die Gleichung beschreibt die Bedingung.
Die Schaltung verbindet Bauteile. Die Tabelle prüft, ob jede Kombination stimmt.

\block{Die Zeichen richtig lesen}
\begin{center}
\begin{tabular}{lll}\toprule
Zeichen&Lesart&Beispiel\\\midrule
$\overline A$&NOT A, A negiert&$\overline0=1$\\
$A\cdot B$ oder $AB$&A AND B&$1\cdot0=0$\\
$A+B$&A OR B&$1+1=1$ in Boolescher Algebra\\
$A\oplus B$&A XOR B&$1\oplus1=0$\\\bottomrule
\end{tabular}
\end{center}
\textbf{Wichtig:} Das Plus in einer Booleschen Gleichung bedeutet OR, keine
gewöhnliche Addition. Rechenreihenfolge: Klammern, NOT, AND, OR.
Ein Überstrich gilt für den gesamten Ausdruck darunter.
\source{DDCA RISC-V, Abschnitt 1.5, S. 17--20}

\newpage
\titlepagepart{02 / NACHSCHLAGEN}{Die Gatter auf einen Blick}
\begin{center}
\begin{tikzpicture}
\node[buffer gate US,draw,scale=1.4] (g0) at (0,0) {};
\draw (g0.input)--++(-0.8,0) node[left] {$A$};
\draw (g0.output)--++(0.8,0) node[right] {$Y$};
\node at (0,-0.85) {\textbf{BUF}\quad $Y=A$};
\node[not gate US,draw,scale=1.4] (g1) at (7,0) {};
\draw (g1.input)--++(-0.8,0) node[left] {$A$};
\draw (g1.output)--++(0.8,0) node[right] {$Y$};
\node at (7,-0.85) {\textbf{NOT}\quad $Y=\overline A$};
\node[and gate US,draw,scale=1.4,logic gate inputs=nn] (g2) at (0,-2) {};
\draw (g2.input 1)--++(-0.8,0) node[left] {$A$};
\draw (g2.input 2)--++(-0.8,0) node[left] {$B$};
\draw (g2.output)--++(0.8,0) node[right] {$Y$};
\node at (0,-2.85) {\textbf{AND}\quad $Y=AB$};
\node[or gate US,draw,scale=1.4,logic gate inputs=nn] (g3) at (7,-2) {};
\draw (g3.input 1)--++(-0.8,0) node[left] {$A$};
\draw (g3.input 2)--++(-0.8,0) node[left] {$B$};
\draw (g3.output)--++(0.8,0) node[right] {$Y$};
\node at (7,-2.85) {\textbf{OR}\quad $Y=A+B$};
\node[nand gate US,draw,scale=1.4,logic gate inputs=nn] (g4) at (0,-4) {};
\draw (g4.input 1)--++(-0.8,0) node[left] {$A$};
\draw (g4.input 2)--++(-0.8,0) node[left] {$B$};
\draw (g4.output)--++(0.8,0) node[right] {$Y$};
\node at (0,-4.85) {\textbf{NAND}\quad $Y=\overline{AB}$};
\node[nor gate US,draw,scale=1.4,logic gate inputs=nn] (g5) at (7,-4) {};
\draw (g5.input 1)--++(-0.8,0) node[left] {$A$};
\draw (g5.input 2)--++(-0.8,0) node[left] {$B$};
\draw (g5.output)--++(0.8,0) node[right] {$Y$};
\node at (7,-4.85) {\textbf{NOR}\quad $Y=\overline{A+B}$};
\node[xor gate US,draw,scale=1.4,logic gate inputs=nn] (g6) at (0,-6) {};
\draw (g6.input 1)--++(-0.8,0) node[left] {$A$};
\draw (g6.input 2)--++(-0.8,0) node[left] {$B$};
\draw (g6.output)--++(0.8,0) node[right] {$Y$};
\node at (0,-6.85) {\textbf{XOR}\quad $Y=A\oplus B$};
\node[xnor gate US,draw,scale=1.4,logic gate inputs=nn] (g7) at (7,-6) {};
\draw (g7.input 1)--++(-0.8,0) node[left] {$A$};
\draw (g7.input 2)--++(-0.8,0) node[left] {$B$};
\draw (g7.output)--++(0.8,0) node[right] {$Y$};
\node at (7,-6.85) {\textbf{XNOR}\quad $Y=\overline{A\oplus B}$};

\end{tikzpicture}
\end{center}
Der kleine Kreis am Ausgang heißt \textbf{Inversionsblase}: Er kehrt das Ergebnis
um. NAND ist AND mit negiertem Ausgang, nicht „beide Eingänge negieren“.

\block{Die vollständige Tabelle für zwei Eingänge}
\begin{center}
\begin{tabular}{cc|rrrrrr}\toprule
$A$&$B$&AND&OR&XOR&NAND&NOR&XNOR\\\midrule
0&0&0&0&0&1&1&1\\
0&1&0&1&1&1&0&0\\
1&0&0&1&1&1&0&0\\
1&1&1&1&0&0&0&1\\\bottomrule
\end{tabular}
\end{center}
\textbf{OR:} mindestens einer ist 1, auch beide dürfen 1 sein.
\textbf{XOR:} bei zwei Eingängen genau einer ist 1.
\textbf{XNOR:} bei zwei Eingängen beide gleich.
NOT tauscht 0 und 1; der Buffer gibt seinen Eingang unverändert weiter.
Ein Buffer ist logisch wie eine Leitung, kann elektrisch aber eine Last treiben.

\block{Mehr als zwei Eingänge}
AND verlangt, dass \emph{alle} Eingänge 1 sind. OR verlangt \emph{mindestens einen}.
XOR mit mehreren Eingängen ist 1 bei einer \textbf{ungeraden Anzahl} von Einsen.
Drei Einsen ergeben daher beim Drei-Eingang-XOR ebenfalls 1.

\block{Verbindung zu Tag 4}
Bei der Addition zweier einzelner Bits $A$ und $B$ gilt:
\[\text{Summenbit}=A\oplus B,\qquad\text{Übertragsbit}=AB.\]
Für $1+1$ entstehen Summenbit 0 und Übertrag 1, also $10_2$.
Das ist ein Halbaddierer; ein zusätzlicher eingehender Übertrag gehört noch nicht dazu.
\source{Abschnitt 1.5, S. 18--20; Halbaddierer als eigene Brücke zu Tag 4}

\newpage
\titlepagepart{03 / BOOLESCHE ALGEBRA}{Gleichungen verstehen und vereinfachen}
Die Variablen haben nur die Werte 0 und 1. Die Grundregeln (Axiome) legen NOT,
AND und OR fest: etwa $\overline0=1$, $\overline1=0$, $0\cdot1=0$ und $1+1=1$.
Daraus folgen die Gesetze, mit denen wir Schaltungen vereinfachen.
\begin{center}
\begin{tabular}{lll}\toprule
Gesetz&AND-Form&OR-Form\\\midrule
Neutrales Element&$A\cdot1=A$&$A+0=A$\\
Dominierendes Element&$A\cdot0=0$&$A+1=1$\\
Idempotenz&$AA=A$&$A+A=A$\\
Komplement&$A\overline A=0$&$A+\overline A=1$\\
Vertauschen&$AB=BA$&$A+B=B+A$\\
Gruppieren&$(AB)C=A(BC)$&$(A+B)+C=A+(B+C)$\\
Verteilen&$A(B+C)=AB+AC$&$A+BC=(A+B)(A+C)$\\
Absorption&$A(A+B)=A$&$A+AB=A$\\
Zusammenfassen&$(A+B)(A+\overline B)=A$&$AB+A\overline B=A$\\\bottomrule
\end{tabular}
\end{center}
Doppelte Negation: $\overline{\overline A}=A$.
\textbf{Dualität:} In einem gültigen Gesetz AND und OR sowie 0 und 1 gleichzeitig
vertauschen ergibt wieder ein gültiges Gesetz. Das ist keine Negation der Variablen.

\block{Warum vereinfachen?}
\[Y=AB+A\overline B=A(B+\overline B)=A\cdot1=A.\]
Die ursprüngliche Beschreibung verwendet mehrere Operationen; logisch genügt
am Ende eine Verbindung von $A$ zu $Y$. Weniger Gatter können Fläche und
Energiebedarf reduzieren. Ob eine konkrete Umsetzung schneller ist, hängt auch
von Gattertypen, Lasten und Leitungen ab.

\block{De Morgan: Negation über eine Verknüpfung ziehen}
\[
\boxed{\overline{AB}=\overline A+\overline B}\qquad
\boxed{\overline{A+B}=\overline A\,\overline B}
\]
Merke: \textbf{Verknüpfung wechseln und jeden Eingang negieren.}
„Nicht beide“ ist „mindestens einer nicht“.
„Keiner von beiden“ ist „beide nicht“.
\textbf{Falsch} wäre $\overline{AB}=\overline A\,\overline B$.

\begin{center}
\begin{tabular}{cc|cc|cc}\toprule
$A$&$B$&$\overline{AB}$&$\overline A+\overline B$&$\overline{A+B}$&$\overline A\,\overline B$\\\midrule
0&0&1&1&1&1\\
0&1&1&1&0&0\\
1&0&1&1&0&0\\
1&1&0&0&0&0\\\bottomrule
\end{tabular}
\end{center}
Jedes Spaltenpaar stimmt in \emph{allen} Zeilen überein: Das beweist hier die
Gleichheit vollständig. Eine einzige geprüfte Kombination wäre kein Beweis.

\block{Weitere Buchregel zum Nachschlagen}
Das Konsensgesetz lautet $AB+\overline A C+BC=AB+\overline A C$.
Der Term $BC$ ist logisch redundant: Ist $B=C=1$, deckt je nach $A$ bereits
$AB$ oder $\overline A C$ diesen Fall ab. Für die ersten NAND-Schaltungen
reichen vor allem doppelte Negation, Idempotenz und De Morgan.
\source{Ergänzend Abschnitt 2.3, S. 58--63; Tabellen 2.1--2.3}

\newpage
\titlepagepart{04 / NAND ALS BAUSTEIN}{Mit einem Gattertyp viele Funktionen bauen}
Wir schreiben kurz $N(A,B)=\overline{AB}$. In den folgenden Zeichnungen ist
jeder beschriftete NAND-Kasten genau ein Zwei-Eingang-NAND; seine Ausgangsnegation
ist bereits im Namen enthalten. Gleich benannte Signale sind verbunden.

\block{NOT aus einem NAND}
Verbinde \emph{beide} Eingänge mit $A$. Dann gilt
$N(A,A)=\overline{AA}=\overline A$.
\begin{center}
\begin{tikzpicture}
\node[gate] (n) at (2,0) {NAND};
\draw[wire] (-1,0) node[left] {$A$}--(0,0);\fill (0,0) circle(1.5pt);
\draw[wire] (0,0)|-($(n.west)+(0,0.25)$);
\draw[wire] (0,0)|-($(n.west)+(0,-0.25)$);
\draw[wire] (n.east)--(4,0) node[right] {$Y=\overline A$};
\end{tikzpicture}
\end{center}

\block{AND aus zwei NANDs}
Erst NAND, dann dessen Ergebnis mit einem weiteren NAND invertieren:
\[T=N(A,B),\qquad Y=N(T,T)=\overline{\overline{AB}}=AB.\]
\begin{center}
\begin{tikzpicture}
\node[gate] (n1) at (0,0) {NAND};\node[gate] (n2) at (4,0) {NAND};
\draw[wire] ($(n1.west)+(0,.25)$)--++(-1,0) node[left] {$A$};
\draw[wire] ($(n1.west)+(0,-.25)$)--++(-1,0) node[left] {$B$};
\draw[wire] (n1.east)--(2,0) node[above] {$T$};\fill (2,0) circle(1.5pt);
\draw[wire] (2,0)|-($(n2.west)+(0,.25)$);
\draw[wire] (2,0)|-($(n2.west)+(0,-.25)$);
\draw[wire] (n2.east)--++(1,0) node[right] {$Y$};
\end{tikzpicture}
\end{center}

\block{OR aus drei NANDs}
Beide Eingänge einzeln negieren und die Ergebnisse NAND-verknüpfen:
\[P=N(A,A),\quad Q=N(B,B),\quad
Y=N(P,Q)=\overline{\overline A\,\overline B}=A+B.\]
\begin{center}
\begin{tikzpicture}
\node[gate] (p) at (0,1) {NAND};\node[gate] (q) at (0,-1) {NAND};
\node[gate] (y) at (4,0) {NAND};
\foreach \g/\v/\h in {p/A/1,q/B/-1}{
\draw[wire] (-2,\h) node[left] {$\v$}--(-1.2,\h);\fill(-1.2,\h)circle(1.5pt);
\draw[wire] (-1.2,\h)|-($(\g.west)+(0,.25)$);
\draw[wire] (-1.2,\h)|-($(\g.west)+(0,-.25)$);}
\draw[wire] (p.east)--(2,1) node[above] {$P$}|-($(y.west)+(0,.25)$);
\draw[wire] (q.east)--(2,-1) node[below] {$Q$}|-($(y.west)+(0,-.25)$);
\draw[wire] (y.east)--++(1,0) node[right] {$Y$};
\end{tikzpicture}
\end{center}
\block{Warum heißt NAND universell?}
Aus NAND können wir NOT, AND und OR bauen. Jede Boolesche Funktion mit endlich
vielen Eingängen lässt sich aus diesen Operationen darstellen, etwa als OR der
AND-Terme für alle Tabellenzeilen mit Ausgang 1. Deshalb genügt NAND allein
für jede solche \emph{kombinatorische} Funktion. Daraus folgt nicht, dass jede
Funktion nur ein NAND benötigt oder dass NAND immer die beste Umsetzung ist.
\source{Ableitungen aus 1.5 und 2.3; Bezug Wahrheitstabelle/Gleichung: 2.2 und 2.4}

\newpage
\titlepagepart{05 / AUFBAU UND NACHWEIS}{XOR und De Morgan in Schaltungen}
\block{XOR aus vier NANDs}
Definiere zuerst Zwischenwerte, statt die ganze Schaltung auf einmal zu lesen:
\[T=N(A,B),\quad P=N(A,T),\quad Q=N(B,T),\quad Y=N(P,Q).\]
\begin{center}
\begin{tikzpicture}
\node[gate] (t) at (0,0) {NAND};
\node[gate] (p) at (4,1) {NAND};\node[gate] (q) at (4,-1) {NAND};
\node[gate] (y) at (8,0) {NAND};
\draw[wire] ($(t.west)+(0,.25)$)--++(-.8,0) node[left] {$A$};
\draw[wire] ($(t.west)+(0,-.25)$)--++(-.8,0) node[left] {$B$};
\draw[wire] (t.east)--(2,0) node[above left] {$T$};\fill(2,0)circle(1.5pt);
\draw[wire] (2,0)|-($(p.west)+(0,-.25)$);
\draw[wire] (2,0)|-($(q.west)+(0,.25)$);
\draw[wire] ($(p.west)+(0,.25)$)--++(-.8,0) node[left] {$A$};
\draw[wire] ($(q.west)+(0,-.25)$)--++(-.8,0) node[left] {$B$};
\draw[wire] (p.east)--(6,1) node[above] {$P$}|-($(y.west)+(0,.25)$);
\draw[wire] (q.east)--(6,-1) node[below] {$Q$}|-($(y.west)+(0,-.25)$);
\draw[wire] (y.east)--++(.8,0) node[right] {$Y$};
\end{tikzpicture}
\end{center}
\begin{center}
\begin{tabular}{cc|cccc}\toprule
$A$&$B$&$T$&$P$&$Q$&$Y$\\\midrule
0&0&1&1&1&0\\
0&1&1&1&0&1\\
1&0&1&0&1&1\\
1&1&0&1&1&0\\\bottomrule
\end{tabular}
\end{center}
Zum Beispiel $A=0,B=1$: $T=1$, $P=1$, $Q=0$, also $Y=N(1,0)=1$.
Die Ergebnisspalte entspricht XOR. Das ist ein vollständiger Funktionsnachweis.

\block{De Morgan als Umzeichnen von Gattern}
\begin{center}
\begin{tikzpicture}
\node[nand gate US,draw,logic gate inputs=nn,scale=1.5] (a) at (0,1) {};
\node[or gate US,draw,logic gate inputs=ii,scale=1.5] (b) at (5,1) {};
\node at (2.5,1) {$\equiv$};
\node[nor gate US,draw,logic gate inputs=nn,scale=1.5] (c) at (0,-1) {};
\node[and gate US,draw,logic gate inputs=ii,scale=1.5] (d) at (5,-1) {};
\node at (2.5,-1) {$\equiv$};
\foreach \g in {a,b,c,d}{
\draw (\g.input 1)--++(-.6,0) node[left] {$A$};
\draw (\g.input 2)--++(-.6,0) node[left] {$B$};
\draw (\g.output)--++(.6,0) node[right] {$Y$};}
\end{tikzpicture}
\end{center}
Oben: NAND entspricht OR mit \emph{beiden} negierten Eingängen.
Unten: NOR entspricht AND mit \emph{beiden} negierten Eingängen.
Eine Ausgangsblase über den Gatterkörper schieben heißt: AND/OR tauschen und
an \emph{jedem} Eingang eine Blase einsetzen. Zwei aufeinanderfolgende Negationen
heben sich logisch auf.

\block{Zwei Nachweise in deiner Übung}
Vergleiche im ersten Versuch $\overline{AB}$ mit $\overline A+\overline B$,
im zweiten $\overline{A+B}$ mit $\overline A\,\overline B$.
Teste jeweils 00, 01, 10 und 11. Beide Ausgänge müssen in jeder Zeile gleich sein.
Ein optionales XOR zwischen den beiden Ausgängen muss dann immer 0 liefern.
\source{De Morgan: 2.3, S. 61--62; Bubble Pushing: 2.5.2; eigene NAND-XOR-Ableitung}

\newpage
\titlepagepart{06 / DIE ELEKTRISCHE SEITE}{Warum 0 und 1 Spannungsbereiche sind}
Ein Draht trägt keine aufgedruckte 0 oder 1, sondern eine kontinuierliche
Spannung. Ein Gatter muss auch leicht gestörte Spannungen richtig lesen.
Deshalb steht nicht nur exakt 0 V für LOW und exakt $V_{DD}$ für HIGH.

\block{Treiber und Empfänger unterscheiden}
Der \textbf{Treiber} erzeugt eine Ausgangsspannung; der \textbf{Empfänger} liest sie.
$V_{DD}$ ist die Versorgungsspannung, GND das Bezugspotential 0 V.
\begin{center}
\begin{tikzpicture}[x=1cm,y=1cm]
\draw[->] (0,0)--(0,4.7) node[above] {$V$};
\draw[->] (6,0)--(6,4.7) node[above] {$V$};
\draw (0,0) rectangle (2,.7);\node at (1,.35) {LOW};
\draw (0,3.7) rectangle (2,4.3);\node at (1,4) {HIGH};
\draw (6,0) rectangle (8,1.3);\node at (7,.6) {LOW};
\draw[pattern=north east lines] (6,1.3) rectangle (8,3);
\draw (6,3) rectangle (8,4.3);\node at (7,3.65) {HIGH};
\node[left] at (0,.7) {$V_{OL}$};\node[left] at (0,3.7) {$V_{OH}$};
\node[right] at (8,1.3) {$V_{IL}$};\node[right] at (8,3) {$V_{IH}$};
\node[left] at (0,4.3) {$V_{DD}$};\node[left] at (0,0) {0};
\draw[dashed] (2,.7)--(5,.7);\draw[dashed] (2,3.7)--(5,3.7);
\draw[dashed] (5,1.3)--(6,1.3);\draw[dashed] (5,3)--(6,3);
\draw[<->] (4.5,.7)--(4.5,1.3) node[midway,left] {$NM_L$};
\draw[<->] (4.5,3)--(4.5,3.7) node[midway,left] {$NM_H$};
\node at (1,-.5) {Treiber-Ausgang};\node at (7,-.5) {Empfänger-Eingang};
\end{tikzpicture}
\end{center}
Schraffiert: kein garantierter Logikwert. Die Zeichnung ist schematisch und
nicht maßstäblich. Sie zeigt keine realen Bauteilkennwerte.
\begin{center}
\begin{tabular}{ll}\toprule
$V_{OL}$&höchste garantierte LOW-Ausgangsspannung\\
$V_{OH}$&niedrigste garantierte HIGH-Ausgangsspannung\\
$V_{IL}$&höchste noch als LOW akzeptierte Eingangsspannung\\
$V_{IH}$&niedrigste schon als HIGH akzeptierte Eingangsspannung\\\bottomrule
\end{tabular}
\end{center}
Zwischen $V_{IL}$ und $V_{IH}$ ist die Interpretation nicht garantiert.
Das ist kein dritter normaler Boolescher Wert.

\block{Störabstand: Wie viel Reserve bleibt?}
\[\boxed{NM_L=V_{IL}-V_{OL}}\qquad\boxed{NM_H=V_{OH}-V_{IH}}\]
Eigene Beispielwerte: $V_{DD}=3{,}3$ V, $V_{OL}=0{,}2$ V,
$V_{IL}=0{,}8$ V, $V_{IH}=2{,}0$ V, $V_{OH}=2{,}8$ V.
Dann sind $NM_L=0{,}6$ V und $NM_H=0{,}8$ V.
Ein Anstieg um 0,5 V auf dem schlechtesten LOW bleibt LOW; ein Abfall um 0,9 V
auf dem schlechtesten HIGH verlässt den garantierten HIGH-Bereich.
\source{Abschnitte 1.6.1--1.6.3, S. 20--22; eigene Grafik nach dem Konzept von Abb. 1.23}

\newpage
\titlepagepart{07 / VOM ANALOGEN ZUM DIGITALEN}{Kennlinie und zuverlässige Verknüpfung}
\block{Ein Inverter ist elektrisch nicht ideal}
Die DC-Übertragungskennlinie zeigt $V_{out}$ als Funktion von $V_{in}$ bei
langsam veränderter Eingabe. Ein idealer Inverter würde abrupt umschalten.
Ein realer Inverter besitzt einen kontinuierlichen Übergang.
\begin{center}
\begin{tikzpicture}[x=1cm,y=1cm]
\begin{scope}
\draw[->](0,0)--(4.5,0) node[right] {$V_{in}$};
\draw[->](0,0)--(0,3.5) node[above] {$V_{out}$};
\draw[thick](0,3)--(2,3)--(2,0)--(4,0);
\node[left]at(0,3){$V_{DD}$};\node[below]at(2,0){$V_{DD}/2$};
\node at(2,-.9){Ideales Modell};
\end{scope}
\begin{scope}[xshift=7cm]
\draw[->](0,0)--(4.5,0) node[right] {$V_{in}$};
\draw[->](0,0)--(0,3.5) node[above] {$V_{out}$};
\draw[thick](0,3).. controls (1.1,3) and (1.5,2.9)..(1.75,2.3)
..controls(2,1.7)and(2.2,.45)..(2.7,.2)..controls(3,.05)and(3.6,0)..(4,0);
\node[left]at(0,3){$V_{DD}$};\node[below]at(4,0){$V_{DD}$};
\node at(2,-.9){Realer Verlauf, schematisch};
\end{scope}
\end{tikzpicture}
\end{center}
Bei kleiner Eingangsspannung ist der Ausgang hoch, bei großer niedrig.
Der reale Übergang muss nicht genau bei $V_{DD}/2$ liegen.
Das Buch bestimmt die Eingangsschwellen an den zwei Stellen, an denen
\[\frac{dV_{out}}{dV_{in}}=-1.\]
Das heißt lokal: Eine kleine Erhöhung am Eingang bewirkt eine gleich große
Absenkung am Ausgang. Die Ableitung musst du heute nicht selbst berechnen;
wichtig ist, dass Logikpegel aus elektrischen Eigenschaften entstehen.

\block{Static Discipline: Gültig hinein, gültig heraus}
Bei gültigen, stabilen Eingangssignalen soll jedes Gatter nach seiner
Einschwingzeit gültige Ausgangssignale liefern. So lassen sich viele Gatter
verbinden, ohne jede Kombination vollständig analog analysieren zu müssen.
Während eines Pegelwechsels durchläuft eine reale Spannung auch Zwischenwerte;
die statische Regel beschreibt nicht das gesamte zeitliche Verhalten.

\block{Passen zwei Gatter elektrisch zusammen?}
Prüfe die garantierten Ausgangspegel des Treibers gegen die Eingangsanforderungen
des Empfängers. Für positive Störreserve benötigen wir
\[V_{OL}<V_{IL},\qquad V_{OH}>V_{IH}.\]
Zusätzlich müssen die zulässigen Eingangsspannungen eingehalten werden.
Ein 5-V-Ausgang darf nicht automatisch an einen 3,3-V-Eingang angeschlossen
werden: Auch ein logisch eindeutiges HIGH kann elektrisch zu hoch sein.
Die konkreten Grenzwerte hängen von der Bauteilfamilie ab.

\block{Was Logisim zeigt und was es vereinfacht}
Eine Logiksimulation hilft dir, Wahrheitstabellen und Verbindungen zu prüfen.
Sie ersetzt keine Untersuchung realer Spannungsgrenzen, Ströme und Lasten.
Ein logisch korrekter NAND-Aufbau ist daher der erste Schritt zur Hardware;
die elektrische Umsetzung kommt später dazu.
\source{Abschnitte 1.6.4--1.6.5, S. 22--24; eigene Kennlinien nach Abb. 1.25}

\newpage
\titlepagepart{08 / DEINE PRAXIS}{Vom Lesen zur geprüften Schaltung}
\block{So liest und prüfst du einen Aufbau}
\begin{enumerate}
\item Benenne Eingänge ($A,B$), Zwischenwerte ($T,P,Q$) und Ausgang ($Y$).
\item Notiere für jedes Gatter seine Gleichung. Gehe von links nach rechts vor.
\item Erstelle die vollständige Wahrheitstabelle mit Zwischenwerten.
\item Baue die Schaltung und teste jede Zeile, nicht nur einen Beispielzustand.
\item Vergleiche mit der gewünschten Funktion und speichere den Aufbau.
\end{enumerate}
Ein Verbindungspunkt zeigt eine elektrische Verbindung. Sich kreuzende Linien
ohne Punkt sind in den hier verwendeten Schaltplänen nicht verbunden.
Ein Ausgang darf mehrere Eingänge speisen; zwei Ausgänge direkt zu verbinden
ist keine Methode, AND oder OR zu bilden.

\block{Praktischer Arbeitsauftrag für Tag 5}
Der Lernplan sieht Logisim-Evolution vor. Lege ein Projekt mit getrennten
Teilschaltungen an. Stelle die Pins auf einzelne Bits und die NANDs auf zwei
Eingänge. Baue die folgenden Funktionen ausschließlich aus NAND:
\begin{center}
\begin{tabular}{lll}\toprule
Funktion&NAND-Anzahl dieses Aufbaus&Prüfung\\\midrule
NOT&1&beide Eingabewerte\\
AND&2&00, 01, 10, 11\\
OR&3&00, 01, 10, 11\\
XOR&4&00, 01, 10, 11\\
NOR (zusätzlicher fünfter Aufbau)&4&00, 01, 10, 11\\\bottomrule
\end{tabular}
\end{center}
Für NOR invertierst du den NAND-basierten OR-Ausgang mit einem weiteren NAND.
Damit erhältst du die fünf NAND-Schaltungen der Erfolgskontrolle.
Der Plan nennt vier Zielfunktionen, fordert aber fünf Schaltungen; NOR ergänzt
hier sinnvoll den fünften Aufbau.

Baue außerdem die beiden De-Morgan-Vergleiche von Seite 5. Hier dürfen zum
Vergleichen auch die normalen NOT-, AND- und OR-Symbole verwendet werden.
Halte für beide Vergleiche alle vier Ausgangspaare fest und speichere dein
Projekt als \code{tag05\_nand.circ} in diesem Ordner. Dieses PDF ist die
Lernunterlage; es enthält noch keine fertig gebaute Logisim-Datei.

\block{Selbstkontrolle mit kurzen Antworten}
\textbf{Warum ist OR nicht XOR?} OR akzeptiert auch zwei Einsen; Zwei-Eingang-XOR nicht.
\par\textbf{Warum funktioniert NOT aus NAND?} $AA=A$, daher $N(A,A)=\overline A$.
\par\textbf{Warum reicht NAND?} NOT, AND und OR lassen sich daraus zusammensetzen.
\par\textbf{Warum alle Tabellenzeilen?} Nur so ist die Funktion für jede Eingabe geprüft.
\par\textbf{Warum Spannungsbereiche?} Reale Signale sind kontinuierlich und gestört.
\par\textbf{Warum nicht nur auswendig lernen?} Du solltest jeweils zwischen Worten,
Gleichung, Tabelle und Schaltung wechseln können.

\block{Lesestellen im vorgesehenen Buch}
\textbf{Harris \& Harris: Digital Design and Computer Architecture, RISC-V Edition.}
Abschnitt 1.5 (Buch S. 17--20) erklärt die Gatter; 1.6 (S. 20--24) erklärt
Pegel, Störabstände, Kennlinien und die statische Regel.
\textbf{Ergänzung zum Lernplan:} Die Booleschen Axiome, Gesetze und De Morgan
stehen in 2.3 (S. 58--63), nicht in 1.6. Für die Schaltungsdarstellung und
Bubble Pushing siehe 2.4--2.5.2 (S. 64--71).
Alle Zeichnungen und didaktischen Beispiele dieser Zusammenfassung wurden neu
für deine Übung erstellt. Sie sind keine übernommenen Buchabbildungen.
\end{document}
